% ============ §8 复习卡 ============
\hdr{§8}{复习卡：一页一张，查表用}

\callout{这四张卡是\textbf{速查台账}，不是新的推导——每一行都能在前面找到出处。考前只看这四张 ＋ \S9，不够时再回对应节。}

\subhdr{8.1　名词：判性一栏，复数一栏}

\begin{minipage}{\linewidth}
\tabtop
\begin{tabular}{@{}>{\raggedright\arraybackslash}p{0.16\linewidth}>{\raggedright\arraybackslash}p{0.30\linewidth}>{\raggedright\arraybackslash}p{0.46\linewidth}@{}}
\thead{要判的事} & \thead{判据} & \thead{出处} \\
\end{tabular}
\tabmid
\begin{tabular}{@{}>{\raggedright\arraybackslash}p{0.16\linewidth}>{\raggedright\arraybackslash}p{0.30\linewidth}>{\raggedright\arraybackslash}p{0.46\linewidth}@{}}
性 & 辅音／-й → 阳；-а／-я → 阴；-о／-е → 中 & \S2.2 \\
性（-ь） & ж／ш／ч／щ＋ь、-ость → 阴；-тель、-арь → 阳；否则查 38 词 & \S2.3 \\
性（特例） & па́па、де́душка、дя́дя、мужчи́на、ю́ноша → 阳 & \S2.2 \\
复数 & 硬词干 ＋ы／а；软词干 ＋и／я & \S2.4 \\
复数（例外） & 24 组特殊复数；只有单数 8 词；只有复数 9 词 & \S2.5 \\
\end{tabular}
\tabbot
\end{minipage}

\subhdr{8.2　形容词、物主代词、指示代词：同一套词尾}

\begin{minipage}{\linewidth}
\tabtop
\begin{tabular}{@{}>{\raggedright\arraybackslash}p{0.14\linewidth}>{\raggedright\arraybackslash}p{0.18\linewidth}>{\raggedright\arraybackslash}p{0.14\linewidth}>{\raggedright\arraybackslash}p{0.16\linewidth}>{\raggedright\arraybackslash}p{0.18\linewidth}@{}}
\thead{性数} & \thead{形容词（硬）} & \thead{软词干} & \thead{物主「我的」} & \thead{指示「这个」} \\
\end{tabular}
\tabmid
\begin{tabular}{@{}>{\raggedright\arraybackslash}p{0.14\linewidth}>{\raggedright\arraybackslash}p{0.18\linewidth}>{\raggedright\arraybackslash}p{0.14\linewidth}>{\raggedright\arraybackslash}p{0.16\linewidth}>{\raggedright\arraybackslash}p{0.18\linewidth}@{}}
阳（原形） & но́вый & си́ний & мой & э́тот \\
阴 & но́вая & си́няя & моя́ & э́та \\
中 & но́вое & си́нее & моё & э́то \\
复 & но́вые & си́ние & мои́ & э́ти \\
\end{tabular}
\tabbot
\end{minipage}

\subhdr{8.3　动词：两组词尾、三时态}

\begin{minipage}{\linewidth}
\tabtop
\begin{tabular}{@{}>{\raggedright\arraybackslash}p{0.14\linewidth}>{\raggedright\arraybackslash}p{0.18\linewidth}>{\raggedright\arraybackslash}p{0.18\linewidth}>{\raggedright\arraybackslash}p{0.16\linewidth}>{\raggedright\arraybackslash}p{0.18\linewidth}@{}}
\thead{人称} & \thead{一组 чита́ть} & \thead{二组 говори́ть} & \thead{过去时} & \thead{将来时} \\
\end{tabular}
\tabmid
\begin{tabular}{@{}>{\raggedright\arraybackslash}p{0.14\linewidth}>{\raggedright\arraybackslash}p{0.18\linewidth}>{\raggedright\arraybackslash}p{0.18\linewidth}>{\raggedright\arraybackslash}p{0.16\linewidth}>{\raggedright\arraybackslash}p{0.18\linewidth}@{}}
я & чита́ю & говорю́ & чита́л／чита́ла & бу́ду чита́ть \\
ты & чита́ешь & говори́шь & чита́л／чита́ла & бу́дешь чита́ть \\
он／она́ & чита́ет & говори́т & чита́л／чита́ла & бу́дет чита́ть \\
мы & чита́ем & говори́м & чита́ли & бу́дем чита́ть \\
вы & чита́ете & говори́те & чита́ли & бу́дете чита́ть \\
они́ & чита́ют & говоря́т & чита́ли & бу́дут чита́ть \\
\end{tabular}
\tabbot
\end{minipage}

\textbf{三个提醒}：第一，词干末尾是 ж／ш／ч／щ 时，单数第一人称写 -у、复数第三人称写 -ут（пишу́／пи́шут）；第二，过去时只看性数，不看人称；第三，将来时＝бу́ду ＋ 不定式，动词不变位。

\subhdr{8.4　数词：一个词三种形式}

\begin{minipage}{\linewidth}
\tabtop
\begin{tabular}{@{}>{\raggedright\arraybackslash}p{0.20\linewidth}>{\raggedright\arraybackslash}p{0.16\linewidth}>{\raggedright\arraybackslash}p{0.16\linewidth}>{\raggedright\arraybackslash}p{0.16\linewidth}>{\raggedright\arraybackslash}p{0.20\linewidth}@{}}
\thead{数字结尾} & \thead{юа́нь} & \thead{час} & \thead{мину́та} & \thead{口诀} \\
\end{tabular}
\tabmid
\begin{tabular}{@{}>{\raggedright\arraybackslash}p{0.20\linewidth}>{\raggedright\arraybackslash}p{0.16\linewidth}>{\raggedright\arraybackslash}p{0.16\linewidth}>{\raggedright\arraybackslash}p{0.16\linewidth}>{\raggedright\arraybackslash}p{0.20\linewidth}@{}}
1、21、31 … & юа́нь & час & мину́та & 一用原形 \\
2、3、4、22--24 … & юа́ня & часа́ & мину́ты & 二三四用单二 \\
5 以上、11--14、0 & юа́ней & часо́в & мину́т & 五以上用复二 \\
\end{tabular}
\tabbot
\end{minipage}

% ============ §9 易错点 ============
\hdr{§9}{易错点：十二条，每条都能指回某个词的某一步}

\begin{enumerate}[label=\textbf{\arabic*.}]
\item \textbf{把 я／ю／е／ё 读成两个音}。在元音后读 [йа]、[йу]，\textbf{在辅音后那个 [й] 消失}（тя 读 [т'а]）。见 \S1.2、\S1.4。
\item \textbf{非重音的 о 读成 [о]}。вода́ 读 [вада́]、молоко́ 读 [малако́]。见 \S1.3。
\item \textbf{把 ш／ж 读软}。ши 是 [шы]、же 是 [жэ]。见 \S1.4。
\item \textbf{в 见谁都清化}。вто́рник 里读 [ф]，但 два、твой 里读 [в]。见 \S1.5。
\item \textbf{把 ИК-3 的中心放在句首}。中心在被问的那一项上，移动中心＝改变问题。见 \S1.6。
\item \textbf{-ь 结尾一律判阴性}。день、портфе́ль、слова́рь 都是阳性。见 \S2.3。
\item \textbf{г／к／х／ж／ш／ч／щ 后写 ы}。*ру́скый、*хоро́шый 都错。见 \S4.3。
\item \textbf{中性指示代词写 *э́тот}。中性要用 \textbf{э́то}（э́то пальто́）。见 \S5.1。
\item \textbf{Я нра́вится}。喜欢的人永远用第三格 \textbf{мне}；动词跟被喜欢的东西一致。见 \S5.3。
\item \textbf{过去时跟人称变}。过去时只看性数：同一个 я，男性 чита́л、女性 чита́ла。见 \S6.3。
\item \textbf{-тся 读成两个音}。у́чится 读 [у́чица]。见 \S6.2。
\item \textbf{21 后面用 юа́ней}。21 用 юа́нь，22--24 用 юа́ня，11--14 用 юа́ней。见 \S7.4。
\end{enumerate}

% ============ §10 缺口清单 ============
\hdr{§10}{缺口清单：课件跳了哪一步，本讲义补在哪}

\begin{itemize}[leftmargin=0em]
\item \gapitem{机制缺口 · У1 p28--p30}{课件把 о 分成「带重音」与「不带重音」两组词，没有说为什么非重音的 о 读 [а]。\S1.3 补：重读音节气流最强，非重读元音发不满，口型向 [а] 靠（реду́кция），并给出「先找重音 → 再看 о → 不在重音上就念 [а]」的操作顺序。}
\item \gapitem{机制缺口 · У2 p9--p12}{课件只有 та—тя—тья 这类对照练习，没有说 я／е／ё／ю 为什么能让前面辅音变软。\S1.2 补：这四个字母本身就是「й ＋ 元音」，在辅音后那个 [й] 脱落，只留下「把前面辅音读软」的作用。}
\item \gapitem{形式改写 · У1 p17}{课件把元音按「硬—软」成对列出（а–я、у–ю、э–е、ы–и、о–ё），只给了配对，没给配对的依据。\S1.2 补：四对是「有没有 й」的差别，ы／и 那一对是例外（不带 й，差别在舌位）。}
\item \gapitem{拼接缺口 · У2 p30 ＋ У3 p2--p3}{「浊辅音站在清辅音前要清化」在 У2 p30，「词末清化」在 У3 p2--p3，两讲隔开，没有一处把两条接起来。\S1.5 合并成一条：浊辅音在\textbf{清辅音前或词末}都要清化，并给出六对清浊表与九组组合表。}
\item \gapitem{压缩缺口 · У3 p2--p3}{两页都是成串词表（нож、шкаф、шарф；тоже、можно、школа……），只给词，没给「为什么六个词正好成对」。\S1.5 补：清浊成对的只有六组，л／м／н／р／й 是响音、不参与清化——这就是表里只有六行的原因。}
\item \gapitem{边界缺口 · У3 p10}{课件在句子上标了 ИК-3 的中心位置，但没说中心可以移动、移动后意思会变。\S1.6 补：Это ³Анна?（问是不是 Anna）与 ³Это Анна?（问这是不是 Anna）是两个问题。}
\item \gapitem{机制缺口 · У4 p11--p13}{课件在 -ь 旁边打了感叹号，只给两栏各 19 个词，没给判断办法。\S2.2 补四条可推规律（ж／ш／ч／щ＋ь、-ость、-тель、-арь），并明确标注「这是归纳，不是课件原文」。}
\item \gapitem{压缩缺口 · У5 p6--p8}{阳性、阴性、中性三张复数表分开列，没有点明它们是同一条规律。\S2.4 合并成一张八行表，并给出总规律「硬词干 ＋ы／а，软词干 ＋и／я」。}
\item \gapitem{机制缺口 · У5 p6}{课件直接写 сестра́ → сёстры，没有解释 е 为什么变 ё。\S2.4 表后补：重音跑到 ё 上（сёстры），这类「重音跑位」没有规则，遇到再记。}
\item \gapitem{机制缺口 · У8 p12--p18}{课件给了 нра́виться ＋ 第三格表，只讲「谁变成什么」，没讲谁做主语、动词跟谁一致。\S5.3 补：主语是「被喜欢的东西」，动词跟它一致（нра́вится／нра́вятся），喜欢的人用第三格。}
\item \gapitem{机制缺口 · 动过 p2--p4}{课件只写「-ТЬ → -Л」，没说为什么过去时不看人称。\S6.3 补：过去时来自古代过去分词，只跟性数一致——所以现在时看人称、过去时看性数。}
\item \gapitem{形式改写 · 数词 p6}{课件把 юа́нь／юа́ня／юа́ней 分成三行列出，没有给判断法。\S7.4 补：名词形式由数字的\textbf{最后一位}决定（1 原形、2--4 单数第二格、5 以上复数第二格），11--14 整体归第三行。}
\end{itemize}

% ============ §11 三句话 ============
\hdr{§11}{三句话总结}

\textbf{第一句：先判性数，再选形式。}——名词判性（\S2.2--2.3）、形容词词尾（\S4.2）、物主代词（\S3.2）、指示代词（\S5.1）、动词过去时（\S6.3），全部是同一件事：\textbf{看清楚那个名词是阳性、阴性、中性还是复数}，后面所有词尾都跟着它走。

\textbf{第二句：一条规律会有好几个「位置」。}——清化有两个位置（词前、词末，\S1.5）、软辅音有三种来源（元音、ь、й，\S1.4）、复数词尾在三个性别上各读一遍（\S2.4）、2／3／4 与 5 以上的名词形式是同一条数词规则（\S7.4）。记住位置，比记住每个位置各是什么更省力。

\textbf{第三句：能回指页码的才算学会。}——本讲义每一节都标了它对应课件哪几页；复习时拿课件翻到那一页，能说出「这一页讲什么、我补了什么」，才算真的读懂了。

% ============ §12 页码对照 ============
\hdr{§12}{与课件的页码对照表}

\begin{minipage}{\linewidth}
\tabtop
\begin{tabular}{@{}>{\raggedright\arraybackslash}p{0.09\linewidth}>{\raggedright\arraybackslash}p{0.48\linewidth}>{\raggedright\arraybackslash}p{0.37\linewidth}@{}}
\thead{简称} & \thead{课件页码范围} & \thead{本讲义对应节} \\
\end{tabular}
\tabmid
\begin{tabular}{@{}>{\raggedright\arraybackslash}p{0.09\linewidth}>{\raggedright\arraybackslash}p{0.48\linewidth}>{\raggedright\arraybackslash}p{0.37\linewidth}@{}}
У1 & p9--p18（字母表、元音成对）、p28--p30（重音与弱化） & \S1.1--\S1.3 \\
У2 & p9--p12（硬软辅音）、p19--p25（Кто／Что э́то）、p30（清化表） & \S1.4、\S1.5、\S1.7 \\
У3 & p2--p3（清浊词表）、p9--p10（调型）、p15--p21（Да／Нет、Где／Когда́） & \S1.5--\S1.7 \\
У4 & p9、p11--p13（性与 -ь 词表）、p17--p19（练习） & \S2.1--\S2.3、\S2.6 \\
У5 & p6--p9（复数与只有单数）、p11--p14（描述练习） & \S2.4--\S2.6 \\
У6 & p6（问候）、p10--p11（物主代词表）、p15--p18（练习）、p26（зову́т） & \S3 \\
У7 & p26--p34（形容词词尾）、p36--p37（颜色）、p41--p42（拼写）、p44--p46（模型） & \S4 \\
У8 & p2--p3（指示代词）、p5--p8（最高级）、p12--p19（нра́виться）、p20--p23（что／потому́ что）、p27（ИК-5） & \S5 \\
动现 & p3--p16（两组变位）、p26--p35（по-ру́сски 与副词） & \S6.1--\S6.2、\S6.6 \\
动过 & p2--p4（-л 型）、p8（时间副词）、p15--p19（星期） & \S6.3、\S6.5、\S7.6 \\
动将 & p6--p8（бу́ду ＋ 不定式）、p12（三时态对话） & \S6.4、\S6.5 \\
数词 & p2--p4（1--1000）、p5--p10（价格）、p11--p17（时间）、p18--p22（序数词）、p23--p25（星期与月份） & \S7 \\
\end{tabular}
\tabbot
\end{minipage}

\clearpage
\subhdr{用完这一册之后：五道自测（答案都在前面）}

\begin{minipage}{\linewidth}
\begin{enumerate}[label=\textbf{\arabic*.}]
\item 注音并说出理由：друг ｜ ло́жка ｜ все ｜ вто́рник ｜ ши ｜ вода́ ｜ молоко́（\S1.3--\S1.5）
\item 判性并写复数：слова́рь ｜ тетра́дь ｜ окно́ ｜ кни́га ｜ го́род（\S2.2--\S2.5）
\item 写 мой 的四个形式，再配 дом ｜ кни́га ｜ окно́ ｜ де́ньги（\S3.2）
\item 写 говори́ть 的六个人称，再写它的过去时四个形式（\S6.1、\S6.3）
\item 用俄语说：2 点 42 分 ｜ 52 元 ｜ 75 元（\S7.4--\S7.5）
\end{enumerate}
\end{minipage}

\flowerrule

\begin{minipage}{\linewidth}
{\small\color{inkgray}
本讲义依据 12 份课件（УРОК 1--8、动词现在／过去／将来时、数词，共 344 页）写成。所有规则要么标注了课件页码，要么给出了完整说明；\S1.2 的 й 双重身份、\S1.3 的弱化机制、\S1.4 的 ш／ж 永硬、\S1.5 的两条清化合并、\S2.2 的 -ь 四条归纳、\S5.3 的 нра́виться 主语、\S6.3 的过去时来历、\S7.4 的三句口诀均为本讲义补出，课件未给。课件中 116 页为无文字层图片页，其中字母表（У1 p9）、元音成对（У1 p17）、清化表（У2 p30）、清浊词表（У3 p3）按图页内容转录，其余图页未转录。\\
按技能规范，讲义页数应落在源课件 0.5--1.0 倍（172--344 页）；本册 35 页，是唯一一处明确违反规范的地方，理由见 \S0.4：语言课课件的 116 页图页与成串朗读操练页不进讲义，规则、词表、例题与练习一条未删。整理日期 2026-10-23。}
\end{minipage}
